\chapter*{2024年\textsf{HDU}「线性代数」期中模拟}

\renewcommand*{\arraystretch}{.9}

\vspace{-.4\baselineskip}\setcounter{page}{1}
\begin{minipage}{.66\textwidth}
  \centering
  \includegraphics[width = .96\linewidth]{border}
  % \begin{tasks}
  %   [ label-width = 5\ccwd, item-indent = 2cm,
  %     after-item-skip = 0pt ](2)
  %   \task[\textbf{开课学院}：] 理学院
  %   \task[\textbf{解析提供}：] \textsf{幻想大白 \& \textsc{X.Y. Xu}}
  %   \task[\textbf{课程编码}：] A0714030
  %   \task[\textbf{\textsf\TeX~排版}：] \textsf{Skyrmion}
  % \end{tasks}
\end{minipage}
\hfill
\begin{minipage}{.15\textwidth}
  \centering
  \begin{tikzpicture}
    \node [ inner sep = 0 ] at (0,0) {
      \qrcode[ height = 4em, version = 4 ]{http://weixin.qq.com/r/sRJiemTEGjllrSAZ90ej}};
    \fill [ white ] (-14em/33,14em/33) rectangle (14em/33,-14em/33);
    \fill [ DarkGreen!60, rounded corners = 4em/33 ]
      (-12em/33,12em/33) rectangle (12em/33,-12em/33);
    \node [ inner sep = 0, white, font = \scriptsize ]
      (weixin) at (0,0) {\faIcon{weixin}};
    \node
      [ inner sep = .5ex, align = center, right,
        xshift = 2em, font = \scriptsize ] at (weixin.east) {未\\央\\学\\社};
  \end{tikzpicture}
\end{minipage}
\begin{minipage}{.15\textwidth}
  \centering
  \begin{tikzpicture}
    \node [ inner sep = 0 ] at (0,0) {
      \qrcode [ height = 4em, version = 4 ]{https://qm.qq.com/q/Oo6vNR8ccU}};
    \fill [ white ] (-14em/33,14em/33) rectangle (14em/33,-14em/33);
    \fill [ MidnightBlue!60, rounded corners = 4em/33 ]
      (-12em/33,12em/33) rectangle (12em/33,-12em/33);
    \node [ inner sep = 0, white, font = \scriptsize ]
      (qq) at (0,0) {\faIcon{qq}};
    \node
      [ inner sep = .5ex, align = center, left,
        xshift = -2em, font = \scriptsize ] at (qq.west) {数\\学\\营};
  \end{tikzpicture}
\end{minipage}

\section{选择题（每题3分，共18分）}
\begin{problem}
  设 $\mathbf A^2 - 2024\mathbf A + 2022\mathbf E = \mathbf O$，
  请问 $(\mathbf A - \mathbf E)^{-1}$为
  \begin{tasks}(4)
    \task[\true] $\mathbf A - 2023\mathbf E$
    \task $\mathbf A + 2024\mathbf E$
    \task $\frac12(\mathbf A-2024\mathbf E)$
    \task $\mathbf A - 2025\mathbf E$
  \end{tasks}
\end{problem}
\scratch
\begin{solution}
  将左式凑写成 $(\mathbf A - \mathbf E)\mathbf B = k\mathbf E$ 的形式，
  得到 $\mathbf B = \mathbf A - 2023\mathbf E$，$k = 1$.
\end{solution}

\begin{problem}
  行列式$
  \begin{vmatrix}
    2x & x & 1 & 2 \\
    1  & x & 1 & -1\\
    3  & 2 & x & 1 \\
    1  & 1 & 1 & x \\
  \end{vmatrix}$对应的多项式中 $x^3$ 的系数为
  \begin{tasks}(4)
    \task $2$
    \task $-2$
    \task $1$
    \task[\true] $-1$
 \end{tasks}
\end{problem}
\scratch
\begin{solution}
  在第一行减去第二行不改变行列式的值，操作完发现所有含有 $x$ 的项都分布在对角线上，
  立即得到含有 $x$ 的项是关于 $x$ 的多项式 $(2x-1)x^3$，那么 $x^3$ 的系数显然为 $-1$.
\end{solution}

\begin{problem}
  设 $D =
  \begin{vmatrix}
    1 & 2  & 3\\
    0 & -1 & 2\\
    2 & 4  & 7
  \end{vmatrix}$，$A_{i j}$ 表示 $D$ 中元素 $a_{ij}$ 的代数余子式，
  则 $A_{13} + A_{23} + A_{33} = $
  \begin{tasks}(4)
    \task[\true] $1$
    \task $-1$
    \task $0$
    \task $5$
  \end{tasks}
\end{problem}
\scratch
\begin{solution}
  根据代数余子式的性质，将 $D$ 的第三列改写为 $(1,1,1)$ 后行列式的值即为所求，
  易得结果为 $1$.
\end{solution}

\begin{problem}
  已知矩阵 $\mathbf A = 
  \begin{bmatrix} 
    1 & 2 & 3 & 4 \\ 
    2 & 4 & 6 & 8 \\ 
    1 & 1 & 1 & 1 \\ 
    0 & 2 & 5 & 7 
  \end{bmatrix}
  $，它的行最简型为
  \begin{tasks}(4)
    \task
    $\begin{bmatrix}
      1 & 0 & 1 & 2 \\ 
      0 & 1 & 2 & 3 \\ 
      0 & 0 & 0 & 0 \\ 
      0 & 0 & 0 & 0 
    \end{bmatrix}$
    \task
    $\begin{bmatrix} 
      1 & 0 & 1 & 2 \\ 
      0 & 1 & 2 & 3 \\ 
      0 & 0 & 1 & 1 \\ 
      0 & 0 & 0 & 0 
    \end{bmatrix}$
    \task[\true]
    $\begin{bmatrix} 
      1 & 0 & 0 & -1 \\ 
      0 & 1 & 0 & 1 \\ 
      0 & 0 & 1 & 1 \\ 
      0 & 0 & 0 & 0 
    \end{bmatrix}$
    \task
    $\begin{bmatrix} 
      1 & 0 & 1 & 2 \\ 
      0 & 1 & 0 & 1 \\ 
      0 & 0 & 0 & 0 \\ 
      0 & 0 & 0 & 0 
    \end{bmatrix}$
  \end{tasks}
\end{problem}
\scratch
\begin{solution}
  行变换得到行阶梯型矩阵
  $\begin{bmatrix}
    1 & 2 & 3 & 4 \\
    0 & 1 & 0 & 1 \\
    0 & 0 & 1 & 1 \\
    0 & 0 & 0 & 0
  \end{bmatrix}$ 后分别用第三行和第二行消除第一行的非 $0$ 元素即可.
\end{solution}

\begin{problem}
  设 $\mathbf N =
  \begin{bmatrix} 
    0 & 1 & 1 & -4 \\ 
    0 & 0 & -5 & 1 \\ 
    0 & 0 & 0 & 4 \\ 
    0 & 0 & 0 & 0 
  \end{bmatrix}$，若 $\mathbf S = \mathbf N^{2024}$， 则 $\mathbf S_{2,4} = $
  \begin{tasks}(4)
    \task $2024$
    \task $-2024$
    \task[\true] $0$
    \task $1$
  \end{tasks}
\end{problem}
\scratch
\begin{solution}
  观察 $\mathbf N$ 的前 $k$ 次方可以轻松发现 $\mathbf N^k$ 的第 $k$ 列为全 $0$，
  因此答案为 $0$.
\end{solution}

\section{填空题（每题3分，共15分）}

\begin{problem}
  设 $\mathbf A^*$ 是三阶方阵 $\mathbf A$ 的伴随矩阵， 若 $|\mathbf A|=2$， 
  则 $||\mathbf A|\mathbf A^*| = $\result{$32$}.
\end{problem}
\scratch[6ex]
\begin{solution}
  $|\mathbf A|\mathbf A^* = |\mathbf A|^2\mathbf A^* = 4\mathbf A^{-1}$，
  因此 $||\mathbf A|\mathbf A^*| = |4\mathbf A^{-1}| = 4^3/|\mathbf A| = 32$.
\end{solution}

\begin{problem}
  已知 $\mathbf A$ 为四阶方阵且 $|\mathbf A| = 2$，
  则 $|(2\mathbf A)^{-1}| = $\result{$1/32$}.
\end{problem}
\scratch[6ex]
\begin{solution}
  由公式 $|\mathbf A^{-1}| = 1/|\mathbf A|$ 和 $|kA| = k^{n}|\mathbf A|$ 得
  $ |(2\mathbf A)^{-1}| = 1/|2\mathbf A| = 1/2^4|\mathbf A|
  = 1/16\times 2 = 1/32$.
\end{solution}

\begin{problem}
  $45321$ 的逆序数是\result{$9$}.
\end{problem}
\scratch[6ex]
\begin{solution}
  $45321$ 中的逆序有
  $(4,3),~ (4,2),~ (4,1),~ (5,3),~ (5,2),~ (5,1),~ (3,2),~ (3,1),~ (2,1)$，
  其逆序数为 $9$.
\end{solution}

\begin{problem}
  设矩阵 $\mathbf A$ 满足 $\mathbf A^2 + \mathbf  A - 4\mathbf E = \mathbf O$，
  则 $(\mathbf A - \mathbf E)^{-1} = $\result{$(\mathbf A + 2\mathbf E)/2$}.
\end{problem}
\scratch[6ex]
\begin{solution}
  $\mathbf A^2 + \mathbf A - 4\mathbf E = \mathbf O \longleftrightarrow
  (\mathbf A - \mathbf E)(\mathbf A + 2\mathbf E)/2 = \mathbf E$，
  由可逆矩阵的定义得：$(\mathbf A - \mathbf E)^{-1} = (\mathbf A + 2\mathbf E)/2$.
\end{solution}

\begin{problem}
  已知齐次线性方程组
  $\begin{cases}
    kx_1 + 2x_2 + x_3  = 0\\
    x_1  + kx_2 + 2x_3 = 0\\
    2x_1 + x_2  + kx_3 = 0
  \end{cases}$ 有非零解，则 $k$ 需要满足的条件为\result{$k = -3$}.
\end{problem}
\scratch[6ex]
\begin{solution}
  令系数矩阵 $\mathbf A =
  \begin{bmatrix}
    k & 2 & 1\\
    1 & k & 2\\
    2 & 1 & k
  \end{bmatrix}$，
  齐次线性方程组有非零解
  $\longleftrightarrow |\mathbf A| = 0 \longleftrightarrow k = -3$.
\end{solution}

\begin{problem}
  写出一个 $n$ 阶方阵 $\mathbf A$ 可逆的充要条件
  \result{存在 $n$ 阶方阵 $\mathbf B$，使 $\mathbf{AB} = \mathbf E$}或
  \result{$\mathbf{AX} = \mathbf O$ 仅有零解}或
  \result{$|\mathbf A|\neq 0$}或
  \result{$\rank(\mathbf A) = n$}.
\end{problem}
\scratch[6ex]

\section{解答题（共64分）}

\begin{problem}[本题3分]
  已知线性变换
  $\begin{cases}
    y_1 = x_1 +  2x_2 - x_3\\
    y_2 = 3x_1 + 4x_2 - 2x_3\\
    y_3 = 5x_1 - 4x_2 + x_3
  \end{cases}$，求 $y_1,~ y_2,~ y_3$ 到 $x_1,~ x_2,~ x_3$ 的线性变换.
\end{problem}
\scratch[10ex]
\begin{solution}
  系数矩阵 $\mathbf A =
  \begin{bmatrix}
    1 & 2 & -1\\
    3 & 4 & -2\\
    5 &-4 & 1
  \end{bmatrix}$. 令 $\mathbf X =
  \begin{bmatrix}
    x_1\\x_2\\x_3
  \end{bmatrix}$，$\mathbf Y = 
  \begin{bmatrix}
    y_1\\y_2\\y_3
  \end{bmatrix}$.
  由 $\mathbf{AX} = \mathbf Y$，$|\mathbf A| \neq 0$，
  所以
  \[
    \mathbf X = \mathbf A^{-1}\mathbf Y =
    \begin{bmatrix}
      -2    & 1 & 0   \\
      -13/2 & 3 & -1/2\\
      -16   & 7 & -1
    \end{bmatrix}
    \begin{bmatrix}
      y_1\\y_2\\y_3
    \end{bmatrix},~
    \begin{cases}
      x_1 = -2y_1    + y_2\\
      x_2 = -13y_1/2 + 3y_2 - y_3/2\\
      x_3 = -16y_1   + 7y_2 - y_3
    \end{cases}
  \]
\end{solution}

\begin{problem}[本题6分]
  求 $x$ 的值：
  $\begin{vmatrix}
    x - 1 & 1 & 2\\
    1     &x-5& 0\\
    2     & 0 & x-5
  \end{vmatrix} = 0$.
\end{problem}
\scratch[8ex]
\begin{solution}
  $D = (x - 5) (x^2 - 6x) = 0$，故 $x = 5$ 或 $x = 0$ 或 $x = 6$.
\end{solution}

\begin{problem}[本题6分]
  若矩阵 $\mathbf A = 
  \begin{bmatrix}
    0 & 1 & 3 \\
    1 & -1 & 0 \\
    -1 & 2 &1 
  \end{bmatrix} $，求 $\mathbf A^{-1}$.
\end{problem}
\scratch[8ex]
\begin{solution}
  由 $\mathbf A^* =
  \begin{bmatrix}
    -1 & 5  & 3 \\
    -1 & 3  & 3\\
    1  & -1 & -1
  \end{bmatrix}$，且 $|\mathbf A| = 2\neq 0$，$\mathbf A$ 可逆，
  故 $\mathbf A^{-1} = \mathbf A^*/|\mathbf A| = \frac12
  \begin{bmatrix}
    -1 & 5  &  3\\
    -3 & 3  &  3\\
    1  & -1 & -1
  \end{bmatrix}$.
\end{solution}

\begin{problem}[本题6分]
  已知 $\mathbf A =
  \begin{bmatrix}
    3 & 1 & 0 & 0\\
    0 & 3 & 0 & 0\\
    0 & 0 & 3 & 9\\
    0 & 0 & 1 & 3
  \end{bmatrix}$，求 $\mathbf A^2$.
\end{problem}
\scratch[8ex]
\begin{solution}
  由 $
  \begin{bmatrix}
    \mathbf A & \mathbf O\\\mathbf O & \mathbf B
  \end{bmatrix}^{n} =
  \begin{bmatrix}
    \mathbf A^n & \mathbf O\\
    \mathbf O   & \mathbf B^n
  \end{bmatrix}$ 得 $\mathbf A^2 =
  \begin{bmatrix}
    9 & 6 & 0 & 0\\
    0 & 9 & 0 & 0\\
    0 & 0 & 18& 54\\
    0 & 0 & 6 & 18
  \end{bmatrix}$.
\end{solution}

\begin{problem}[本题6分]
  设 $D = 
  \begin{bmatrix}
    1  & 1 & 2 & -1\\
    -2 & 3 & 4 & 1 \\
    3  & 4 & 1 & 2 \\
    -4 & 2 & 0 & 6
  \end{bmatrix}$，
  求下列代数余子式的和.
  \begin{tasks}[label = \arabic*.](2)
    \task $A_{12} - 2A_{22} +
          3A_{32} - 4A_{42}$.
    \task $A_{31} + 2A_{32} + A_{34}$.
  \end{tasks}
\end{problem}
\scratch[8ex]
\begin{solution}\leavevmode
  \begin{enumerate}
    \item 将第二列替换 为 $1,~ -2,~ 3,~ -4$，得到矩阵 $\mathbf B =
    \begin{bmatrix}
      1  &  1 & 2 & -1\\
      -2 & -2 & 4 & 1 \\
      3  &  3 & 1 & 2 \\
      -4 & -4 & 0 & 6
    \end{bmatrix}$，其一二列相同，所以 $|\mathbf B| = 0$
    \item 同理，将第三行替换为 $1,~ 2,~ 0,~ 1$，得 $|\mathbf C| =
    \begin{vmatrix}
      1  & 1 & -2 & -4\\
      -2 & 3 &  4 & 1 \\
      1  & 2 &  0 & 1 \\
      -4 & 2 &  0 & 6
    \end{vmatrix} = -40$.
  \end{enumerate}
\end{solution}

\begin{problem}[本题8分]
  已知 $\mathbf A = 
  \begin{bmatrix}
    2 & 0 & 1\\
    0 & 3 & 0\\
    2 & 0 & 2
  \end{bmatrix},~
  \mathbf B = 
  \begin{bmatrix}
    1 & 0 & 0\\
    0 &-1 & 0\\
    0 & 0 & 0
  \end{bmatrix}$. 若 $\mathbf X$ 满足
  $\mathbf{AX} + 2\mathbf B = \mathbf{BA} + 2\mathbf X$，求 $\mathbf X^4$.
\end{problem}
\scratch[14ex]
\begin{solution}\leavevmode
  \begin{itemize}
    \item $\mathbf{AX} - 2\mathbf X = \mathbf{BA} - \mathbf{AB}
    \longleftrightarrow
      (\mathbf A - 2\mathbf E)\mathbf X = \mathbf B(\mathbf A - 2\mathbf E)$.
    \item 由于 $\mathbf A - 2\mathbf E =
    \begin{bmatrix}
      0 & 0 & 1\\
      0 & 1 & 0\\
      2 & 0 & 0
    \end{bmatrix}$，$|\mathbf A - 2\mathbf E| \neq 0$，
    $\mathbf A - 2\mathbf E$ 可逆，故
    $ \mathbf X^4 =
      (\mathbf A - 2\mathbf E)^{-1}\mathbf B^4(\mathbf A - 2\mathbf E)=
    \begin{bmatrix}
      0 & 0 & 0\\
      0 & 1 & 0\\
      0 & 0 & 1
    \end{bmatrix}$.
  \end{itemize}
\end{solution}

\begin{problem}[本题8分]
  若 $|\mathbf A| = 1,~ |\mathbf B| = 2,~ |\mathbf A + \mathbf B| = 2$，
  求 $|(\mathbf A^{-1} + \mathbf B^{-1})^{-1}|$.
\end{problem}
\scratch[14ex]
\begin{solution}
  $ |(\mathbf A^{-1} + \mathbf B^{-1})^{-1}|
  = \frac1{|\mathbf A^{-1} + \mathbf B^{-1}|}
  = \frac1{|\mathbf A^{-1}||\mathbf E + \mathbf{AB}^{-1}|}
  = \frac1{|\mathbf A^{-1}||\mathbf A + \mathbf B||\mathbf B^{-1|}}
  = \frac{|\mathbf A||\mathbf B|}{|\mathbf A + \mathbf B|} = 1$.
\end{solution}

\begin{problem}[本题8分]
  已知矩阵 $\mathbf A =
  \begin{bmatrix}
    3 & 1 & 1 & 1\\
    1 & 4 & 3 & 4\\
    1 & 4 & 11&16\\
    1 & 8 & 2 &66
  \end{bmatrix}$，若 $\mathbf X$ 满足
  $\mathbf{AX} + 4\mathbf E = \mathbf A^2 + 2\mathbf X$，求 $\mathbf X$.
\end{problem}
\begin{remark}
  $\mathbf{AC} = \mathbf{BC}\nRightarrow \mathbf A = \mathbf B$.
\end{remark}
\scratch[12ex]
\begin{solution}
  $\mathbf{AX} - 4\mathbf E = \mathbf A^2 + 2\mathbf X\longleftrightarrow
    (\mathbf A - 2\mathbf E)(\mathbf A + 2\mathbf E) =
    (\mathbf A - 2\mathbf E)\mathbf X \longleftrightarrow
    (\mathbf A - 2\mathbf E)(\mathbf A + 2\mathbf E - \mathbf X) = \mathbf O$.
  由于 $|\mathbf A - 2\mathbf E| \neq 0$，
  故齐次线性方程组 $\mathbf{AX} = \mathbf O$ 仅有零解，即
  \[\mathbf A + 2\mathbf E - \mathbf X = \mathbf O\longleftrightarrow
  \mathbf X = \mathbf A + 2\mathbf E =
  \begin{bmatrix}
    5 & 1 & 1 & 1\\
    1 & 6 & 3 & 4\\
    1 & 4 &13 &16\\
    1 & 8 & 2 &68
  \end{bmatrix}\]
\end{solution}

\begin{problem}[本题10分]
  设线性方程组
  $\begin{cases}
    x_1 + x_2 + 2x_3 + 3x_4 = 1\\
    x_1 + 3x_2+ 6x_3 + x_4  = 3\\
    3x_1- x_2 - ax_3 + 15x_4= 3\\
    x_1 - 5x_2- 10x_3+ 12x_4= b
  \end{cases}$，问 $a$ 与 $b$ 各取何值时，线性方程组有唯一解？
  无解？无穷多解？并求出无穷多解时的通解.
\end{problem}
\scratch[18ex]
\begin{solution}
  \[
    \overline{\mathbf A} =
    \begin{bmatrix}
      1 & 1 & 2 & 3 & 1\\
      1 & 3 & 6 & 1 & 3\\
      3 &-1 &-a &15 & 3\\
      1 &-5 &-10&12 & b\\
    \end{bmatrix} =
    \begin{bmatrix}
      1 & 0 & 0 & 0 & -4(b+5)/3\\
      0 & 1 & 0 & 0 & (b+8)/3\\
      0 & 0 &2-a& 0 & -2(b-1)/3\\
      0 & 0 & 0 & 1 & (b+5)/3\\
    \end{bmatrix}
  \]
  \begin{itemize}
    \item 当 $a = 2$ 且 $ b\neq 1$ 时，
    $\rank(\mathbf A) = 3,~ \rank(\overline{\mathbf A}) = 4$.
    因为 $\rank(\mathbf A) < \rank(\overline{\mathbf A})$，故此方程组无解.
    \item 当 $a \neq 2$ 时，
    $\rank(\mathbf A) = \rank(\overline{\mathbf A}) = 4$，方程组有唯一解.
    \item 当 $a = 2$ 且 $b = 1$ 时，
    $\rank(\mathbf A) = \rank(\overline{\mathbf A}) = 3$，方程组有无穷多解.
  \end{itemize}

  此时 $\overline{\mathbf A}=
  \begin{bmatrix}
    1 & 0 & 0 & 0 &-8\\
    0 & 1 & 0 & 0 & 3\\
    0 & 0 & 0 & 0 & 0\\
    0 & 0 & 0 & 1 & 2\\
  \end{bmatrix}$. 取基础解系 $\bm\epsilon =
  \begin{bmatrix}
    0\\2\\1\\0
  \end{bmatrix}$，
  通解 $\bm x = k\bm\epsilon +
  \begin{bmatrix}
    -8\\3\\0\\2
  \end{bmatrix}$，其中 $k$ 为任意常数.
\end{solution}

\begin{problem}[本题6分]
  设 $\mathbf A$ 为 $n$ 阶非零矩阵，若 $\mathbf A^3 = \mathbf O$，
  试证明：矩阵 $\mathbf A + \mathbf E,~ \mathbf A - \mathbf E$ 均可逆.
\end{problem}
\scratch[10ex]
\begin{solution}
  \begin{proof}
    采用立方和公式 $\mathbf A^3 = \mathbf O\longleftrightarrow
    \mathbf A^3 + \mathbf E = \mathbf E\longleftrightarrow
    (\mathbf A + \mathbf E)(\mathbf A^2 - \mathbf A + \mathbf E) = \mathbf E$.
    由可逆矩阵的定义得：$\mathbf A + \mathbf E$ 可逆.
    同理，在等式左右同时减去 $\mathbf E$，
    使用立方差公式，可得 $\mathbf A - \mathbf E$ 可逆.
  \end{proof}
\end{solution}